\documentclass[11pt]{article}
\usepackage{mathblatt}


\begin{document}
\pruefheftstil{Prozent}{Prüfungsheft P10}
\noindent{\Large\bfseries Prozent}\hfill{\small\color{mbgrau}Prüfungsheft P10}\par\medskip
\begin{pfuebersicht}
\pfuezeile{sprozentundanteilumwa}{Prozent und Anteil umwandeln}{Anteil, Prozentsatz, Bruch, Prozent, Hundertstel}
\pfuezeile{sprozentwert}{Prozentwert}{Grundwert, Prozentsatz, Rabatt, Differenz, Anzahl}
\pfuezeile{sprozentsatz}{Prozentsatz}{Anteil, Grundwert, Milliarden, Kreisdiagramm, Mittelpunktswinkel}
\pfuezeile{sgrundwert}{Grundwert}{Prozentsatz, Rabatt, Gewinn}
\pfuezeile{serhhungundvernderung}{Erhöhung und Veränderung in Prozent}{Grundwert, Prozentsatz, Spannweite, prozentuale Zunahme, Preiserhöhung}
\pfuezeile{saussagenprfen}{Aussagen prüfen}{Anteil, Steigung in Prozent, Prozentsatz, Behauptung prüfen, Bruch und Prozent}
\pfuezeile{sprozentauseinerberec}{Prozent aus einer berechneten Fläche}{Verschnitt, Kreisfläche, Prozentsatz, Abrunden, Modellieren}
\pfuezeile{szinsenundzinssatz}{Zinsen und Zinssatz}{Zinsen, Zinssatz, Kapital, Prozentsatz, Nachweis}
\pfuezeile{szinseszinsundguthabe}{Zinseszins und Guthabentabelle}{Zinsen, Guthaben, Einzahlung, Tabelle ergänzen, Zinseszins}
\multicolumn{3}{@{}l}{\hspace{1.4em}Prüfstein} & S.\,\pageref{pruefstein}\\
\end{pfuebersicht}
\pfabschnitt{Prozent}
\pfstufe[sprozentundanteilumwa]{Prozent und Anteil umwandeln}
\kommtdran{in 2 von 13 Jahren (2020, 2022) \textperiodcentered{} zuletzt 2022 \textperiodcentered{} 1 BE}
\begin{pfaufgabe}{1.}{1 BE \textperiodcentered{} P10 2022 · 1f}
Ein Fünftel aller Jugendlichen erhält kein Taschengeld. Kreuzen Sie den passenden Prozentsatz an:
\par\noindent\hspace*{1.6em}$\square$ 5 \% \quad $\square$ 20 \% \quad $\square$ 25 \% \quad $\square$ 50 \%\par
\fusshilfe{1: 20 \%, Tipp: Prozentsatz}
\end{pfaufgabe}
\pfweiterekopf
\pfpaar{\pfkopfzeile{2.}{eigene Aufgabe}
Schreibe den Bruch $\tfrac{11}{20}$ in Prozent.
\par\noindent\hfill \leerfeld[\%]
\fusshilfe{2: $55\,\%$}}{\pfkopfzeile{3.}{eigene Aufgabe}
Jeder hundertste Besucher bekommt einen Preis. Wie viel Prozent der Besucher bekommen einen Preis?
\par\noindent\hfill \leerfeld[\%]
\fusshilfe{3: $1\,\%$}}
{\small\begin{pfaufgabe}{4.}{1 BE \textperiodcentered{} P10 2020 · 1a}
Ein Förster sagt: „Bei mir wachsen nur 4 \% der gepflanzten Bäume nicht an.“ Kreuzen Sie an, welche Aussage dasselbe bedeutet:
\pfkreuzzeile{Von 10 Bäumen wachsen 4 nicht an.}\pfkreuzzeile{Jeder vierte Baum wächst nicht an.}\pfkreuzzeile{Von 100 Bäumen wachsen 4 nicht an}
\fusshilfe{4: 4 von 100 Bäumen wachsen nicht an, Tipp: Hundertstel}
\end{pfaufgabe}}
\pfpaar{\pfkopfzeile{5.}{eigene Aufgabe nach P10 2020}
$5\,\%$ der Pakete kommen zu spät. Kreuze die Aussage an, die dazu passt.\\ \kreuz{$5$ von $10$ Paketen kommen zu spät.}\\ \kreuz{Jedes 5. Paket kommt zu spät.}\\ \kreuz{$5$ von $100$ Paketen kommen zu spät.}
\fusshilfe{5: $5$ von $100$ Paketen kommen zu spät}}{}
\pfstufe[sprozentwert]{Prozentwert}
\kommtdran{in 6 von 13 Jahren (2014, 2015, 2017, 2019, 2021, 2026) \textperiodcentered{} zuletzt 2026 \textperiodcentered{} 1--3 BE}
\begin{pfaufgabe}{6.}{1 BE \textperiodcentered{} P10 2026 FOR · 1a}
Geben Sie 30 \% von 70 € an.
\rechenplatz{3}
\fusshilfe{6: 21 €, Tipp: Prozentsatz}
\end{pfaufgabe}
\pfweiterekopf
\pfpaar{\pfkopfzeile{7.}{eigene Aufgabe}
Berechne $20\,\%$ von $4{,}50\,€$.
\par\noindent\hfill \leerfeld[€]
\fusshilfe{7: $0{,}90\,€$}}{\pfkopfzeile{8.}{1 BE \textperiodcentered{} P10 2021 · 1c}
Ein Fahrrad kostet 550 €. Wer bar bezahlt, bekommt 20 \% Rabatt. Kreuzen Sie an, wie viel Euro man dabei spart:
\par\noindent\hspace*{1.6em}$\square$ 100 € \quad $\square$ 440 € \quad $\square$ 110 € \quad $\square$ 660 €\par
\fusshilfe{8: 110 €, Tipp: Rabatt}}
\pfpaar{\pfkopfzeile{9.}{2 BE \textperiodcentered{} P10 2019 · 5a}
In einer Studie wurden dieselben 1 200 Jugendlichen über mehrere Jahre befragt. 2013 hatten 72 \% von ihnen ein Smartphone, 2016 waren es 95 \%. Berechnen Sie, um wie viele Jugendliche die Zahl der Smartphone-Besitzer in dieser Zeit gestiegen ist.
\fusshilfe{9: 276 Jugendliche, Tipp: Differenz}}{\pfkopfzeile{10.}{1 BE \textperiodcentered{} P10 2017 · 1b}
Für eine Bohrmaschine zum Preis von 120,00 € gibt es an der Kasse 20 \% Nachlass. Geben Sie an, wie viel Euro der Nachlass beträgt.
\fusshilfe{10: 24 €, Tipp: Rabatt}}
\pfpaar{\pfkopfzeile{11.}{1 BE \textperiodcentered{} P10 2014 · 1a}
Bestimmen Sie 13 \% von 50 €.
\fusshilfe{11: 6,50 €, Tipp: Grundwert}}{\pfkopfzeile{12.}{eigene Aufgabe nach P10 2021}
Ein Fernseher kostet $480\,€$. Bei Barzahlung gibt es $15\,\%$ Rabatt. Kreuze an, wie viel Euro man spart.\\ \kreuz{$32\,€$}\\ \kreuz{$408\,€$}\\ \kreuz{$72\,€$}\\ \kreuz{$552\,€$}
\fusshilfe{12: $72\,€$}}
\pfpaar{\pfkopfzeile{13.}{eigene Aufgabe}
Dieselbe Jacke kostet in Laden A $95\,€$ mit $20\,\%$ Rabatt. In Laden B kostet sie $86\,€$ mit $10\,\%$ Rabatt. In welchem Laden ist die Jacke billiger?
\fusshilfe{13: $1{,}40\,€$, in Laden A ist die Jacke billiger}}{}
\pfstufe[sprozentsatz]{Prozentsatz}
\kommtdran{in 3 von 13 Jahren (2015, 2018, 2023) \textperiodcentered{} zuletzt 2023 \textperiodcentered{} 1--4 BE}
\begin{pfaufgabe}{14.}{2 BE \textperiodcentered{} P10 2023 · 6a}
Im Jahr 2020 gab es auf der Erde rund 7,8 Milliarden Menschen. Die Tabelle teilt sie nach dem Alter auf: Kinder (0 bis 14 Jahre) 1,95 Milliarden, Jugendliche (15 bis 24 Jahre) 1,25 Milliarden, mittleres Alter (25 bis 64 Jahre) 3,9 Milliarden, höheres Alter (ab 65 Jahre) 0,7 Milliarden. Berechnen Sie den Anteil der Jugendlichen an der Weltbevölkerung im Jahr 2020 in Prozent.
\rechenplatz{3}
\fusshilfe{14: $\approx$ 16,0 \%, Tipp: Anteil}
\end{pfaufgabe}
\pfweiterekopf
\pfpaar{\pfkopfzeile{15.}{eigene Aufgabe}
$17$ von $24$ Kindern waren im Schwimmbad. Wie viel Prozent der Kinder waren im Schwimmbad? Runde auf eine Stelle nach dem Komma.
\par\noindent\hfill \leerfeld[\%]
\fusshilfe{15: $\approx 70{,}8\,\%$}}{\pfkopfzeile{16.}{1 BE \textperiodcentered{} P10 2018 · 7a}
Für ein Fest backt Eva 16 Pfannkuchen: 14 füllt sie mit Marmelade, 2 mit Senf. Geben Sie an, wie viel Prozent der Pfannkuchen eine Senffüllung haben.
\fusshilfe{16: 12,5 \%, Tipp: Anteil}}
\pfpaar{\pfkopfzeile{17.}{eigene Aufgabe nach P10 2018}
Bei einer Tombola gibt es $21$ Nieten und $3$ Hauptgewinne. Wie viel Prozent der Lose sind Hauptgewinne?
\par\noindent\hfill \leerfeld[\%]
\fusshilfe{17: $12{,}5\,\%$}}{}
{\small\begin{pfaufgabe}{18.}{4 BE \textperiodcentered{} P10 2015 · 7c}
In einer Saison der 2. Fußball-Bundesliga wurden 83 Elfmeter gepfiffen; 71 davon führten zu einem Tor. Berechnen Sie, wie viel Prozent der Elfmeter nicht verwandelt wurden. Zeichnen Sie den Anteil als Sektor in den vorbereiteten Kreis ein. Geben Sie an, wie groß der Mittelpunktswinkel dieses Sektors ist.
\par\smallskip\noindent \kreisleer[2]\par
\fusshilfe{18: $\approx$ 14,5 \%; Sektor $\approx$ 52°; $\approx$ 52°, Tipp: Kreisdiagramm}
\end{pfaufgabe}}
\pfpaar{\pfkopfzeile{19.}{eigene Aufgabe}
Ein Handyvertrag hat $25$ GB im Monat. Nach $22$ Tagen sind $16{,}5$ GB verbraucht. Der Monat hat noch $8$ Tage. Wie viel Prozent der $25$ GB sind noch übrig? Reicht der Rest bei gleichem Verbrauch bis zum Monatsende?
\fusshilfe{19: Ja}}{}
\pfstufe[sgrundwert]{Grundwert}
\kommtdran{in 2 von 13 Jahren (2023, 2025) \textperiodcentered{} zuletzt 2025 \textperiodcentered{} 1 BE}
\begin{pfaufgabe}{20.}{1 BE \textperiodcentered{} P10 2025 · 1a}
Ein T-Shirt ist um 6 € billiger geworden. Dieser Nachlass entspricht 20 \% des ursprünglichen Preises. Geben Sie den ursprünglichen Preis an.
\rechenplatz{3}
\fusshilfe{20: 30 €, Tipp: Prozentsatz}
\end{pfaufgabe}
\pfweiterekopf
\pfpaar{\pfkopfzeile{21.}{eigene Aufgabe}
$36\,\%$ einer Menge sind $81$ kg. Wie viel kg ist die ganze Menge?
\par\noindent\hfill \leerfeld[kg]
\fusshilfe{21: $225$ kg}}{\pfkopfzeile{22.}{1 BE \textperiodcentered{} P10 2023 · 1b}
Mia gewinnt bei einem Wettbewerb Geld. Sie gibt 25 \% des Gewinns an ihre Eltern ab; das sind 200 €. Geben Sie an, wie hoch der gesamte Gewinn war.
\fusshilfe{22: 800 €, Tipp: Prozentsatz}}
\pfpaar{\pfkopfzeile{23.}{eigene Aufgabe nach P10 2025}
Beim Kauf einer Mütze spart Mira $4\,€$, das sind $25\,\%$ des alten Preises. Wie viel kostete die Mütze vorher?
\par\noindent\hfill \leerfeld[€]
\fusshilfe{23: $16\,€$}}{\pfkopfzeile{24.}{eigene Aufgabe}
Tims Handy zeigt $28\,\%$ Akku. Laut Anzeige reicht das noch für $7$ Stunden Musik. Wie lange hält ein voller Akku? Reicht ein voller Akku für eine Busfahrt von $22$ Stunden?
\fusshilfe{24: Ja}}
\pfstufe[serhhungundvernderung]{Erhöhung und Veränderung in Prozent}
\kommtdran{in 7 von 13 Jahren (2015, 2016, 2020, 2021, 2022, 2024, 2026) \textperiodcentered{} zuletzt 2026 \textperiodcentered{} 1--5 BE}
\begin{pfaufgabe}{25.}{1 BE \textperiodcentered{} P10 2024 · 1e}
Ein Schwimmbad verlangt bisher 3,50 € Eintritt. Der Preis steigt um 20 \%. Geben Sie den neuen Eintrittspreis in Euro an.
\rechenplatz{3}
\fusshilfe{25: 4,20 €, Tipp: Preiserhöhung}
\end{pfaufgabe}
\pfweiterekopf
\pfpaar{\pfkopfzeile{26.}{eigene Aufgabe}
In einer AG steigt die Zahl der Kinder von $40$ auf $50$. Um wie viel Prozent ist die Zahl gestiegen?
\par\noindent\hfill \leerfeld[\%]
\fusshilfe{26: $25\,\%$}}{}
{\small\begin{pfaufgabe}{27.}{2 BE \textperiodcentered{} P10 2022 · 4b}
Das Säulendiagramm zeigt, wie viele von jeweils 1200 befragten Jugendlichen in den Jahren 2017 bis 2021 mehrmals pro Woche ein Buch lesen. Ermitteln Sie, um wie viel Prozent diese Zahl von 2020 auf 2021 zurückgegangen ist.
\par\smallskip\noindent \begin{tikzpicture}\begin{axis}[ybar,width=8.5cm,height=4.6cm,ymin=390,ymax=508.75,symbolic x coords={2017,2018,2019,2020,2021},xtick=data,enlarge x limits=0.08,bar width=6mm,ylabel={Anzahl},ylabel style={font=\small},tick label style={font=\small},nodes near coords,every node near coord/.append style={font=\scriptsize,/pgf/number format/.cd,use comma,1000 sep={}},ymajorgrids,grid style={mbgitter}]\addplot[fill=mbkasten,draw=black] coordinates {(2017,468.0) (2018,432.0) (2019,485.0) (2020,470.0) (2021,400.0)};\end{axis}\end{tikzpicture}\par
\fusshilfe{27: $\approx$ 14,9 \%, Tipp: prozentuale Abnahme}
\end{pfaufgabe}}
{\small\begin{pfaufgabe}{28.}{5 BE \textperiodcentered{} P10 2021 · 5a}
Das Säulendiagramm gibt an, wie viele Container (in Millionen) in den Jahren 2009 bis 2017 im Hamburger Hafen be- und entladen wurden. Berechnen Sie die Spannweite und den Durchschnitt (das arithmetische Mittel) dieser Containerzahlen. Berechnen Sie, um wie viel Prozent die Containerzahl von 2010 auf 2011 gestiegen ist.
\par\smallskip\noindent \begin{tikzpicture}\begin{axis}[ybar,width=13.7cm,height=4.6cm,ymin=0,ymax=11.64,symbolic x coords={2009,2010,2011,2012,2013,2014,2015,2016,2017},xtick=data,enlarge x limits=0.08,bar width=6mm,ylabel={Container in Millionen},ylabel style={font=\small},tick label style={font=\small},nodes near coords,every node near coord/.append style={font=\scriptsize,/pgf/number format/.cd,use comma,1000 sep={}},ymajorgrids,grid style={mbgitter}]\addplot[fill=mbkasten,draw=black] coordinates {(2009,7.0) (2010,7.9) (2011,9.0) (2012,8.9) (2013,9.3) (2014,9.7) (2015,8.8) (2016,8.9) (2017,8.8)};\end{axis}\end{tikzpicture}\par
\fusshilfe{28: 2,7 Mio.; 8,7 Mio.; $\approx$ 13,9 \%, Tipp: Spannweite}
\end{pfaufgabe}}
\pfpaar{\pfkopfzeile{29.}{eigene Aufgabe nach P10 2026}
Ein Liter Milch kostete $1{,}15\,€$, jetzt kostet er $1{,}29\,€$. Um wie viel Prozent ist der Preis gestiegen? Runde auf eine Stelle nach dem Komma.
\fusshilfe{29: $\approx 12{,}2\,\%$}}{\pfkopfzeile{30.}{2 BE \textperiodcentered{} P10 2026 FOR · 3c}
An einer Tankstelle wurde eine Woche lang täglich um 18 Uhr der Benzinpreis notiert (€ pro Liter): Mo 1,77, Di 1,78, Mi 1,77, Do 1,79, Fr 1,84, Sa 1,82, So 1,85. Berechnen Sie die prozentuale Preissteigerung von Donnerstag auf Freitag.
\fusshilfe{30: $\approx$ 2,8 \%, Tipp: prozentuale Veränderung}}
{\small\begin{pfaufgabe}{31.}{2 BE \textperiodcentered{} P10 2020 · 2d}
Das Säulendiagramm zeigt die Besucherzahlen eines Freibads an den Tagen der ersten Ferienwoche. Die beiden folgenden Aussagen sind richtig: „Die Besucherzahlen der ersten Ferienwoche haben eine Spannweite von 150.“ „Am Sonntag gab es ein Drittel mehr Besucher als am Mittwoch.“ Weisen Sie beide Aussagen rechnerisch nach.
\par\smallskip\noindent \begin{tikzpicture}\begin{axis}[ybar,width=11.1cm,height=4.6cm,ymin=0,ymax=336.00,symbolic x coords={Mo,Di,Mi,Do,Fr,Sa,So},xtick=data,enlarge x limits=0.08,bar width=6mm,ylabel={Besucher je Wochentag},ylabel style={font=\small},tick label style={font=\small},nodes near coords,every node near coord/.append style={font=\scriptsize,/pgf/number format/.cd,use comma,1000 sep={}},ymajorgrids,grid style={mbgitter}]\addplot[fill=mbkasten,draw=black] coordinates {(Mo,170.0) (Di,130.0) (Mi,210.0) (Do,180.0) (Fr,190.0) (Sa,240.0) (So,280.0)};\end{axis}\end{tikzpicture}\par
\fusshilfe{31: 280 $-$ 130 = 150; 210 + 210 : 3 = 280, Tipp: Spannweite}
\end{pfaufgabe}}
{\small\begin{pfaufgabe}{32.}{3 BE \textperiodcentered{} P10 2020 · 4c}
Nach einer Studie lag der Wasserverbrauch 2019 bei durchschnittlich 54 Tausend Litern pro Person und Jahr; er soll jedes Jahr um 3 \% steigen. Diese Vorhersage beschreibt die Gleichung $y = 54\cdot 1{,}03^{x}$ ($x$: Jahre ab 2019, $y$: Verbrauch pro Person in Tausend Litern). Eine zweite Studie erwartet, dass der Verbrauch pro Person von 2019 bis 2033 insgesamt um etwa 40 \% zunimmt. Untersuchen Sie, ob beide Studien für 2033 denselben Verbrauch vorhersagen.
\fusshilfe{32: nein, Studie 1 $\approx$ 81,7 gegen Studie 2 = 75,6 (Tausend Liter), Tipp: Exponentialfunktion}
\end{pfaufgabe}}
{\small\begin{pfaufgabe}{33.}{2 BE \textperiodcentered{} P10 2016 · 2c}
Das Säulendiagramm zeigt, wie viele Besucher ein Filmpark in den Jahren 2007 bis 2015 hatte. Berechnen Sie, um wie viel Prozent die Besucherzahl von 2009 auf 2010 zugenommen hat.
\par\smallskip\noindent \begin{tikzpicture}\begin{axis}[ybar,width=13.7cm,height=4.6cm,ymin=0,ymax=720.00,symbolic x coords={2007,2008,2009,2010,2011,2012,2013,2014,2015},xtick=data,enlarge x limits=0.08,bar width=6mm,ylabel={Besucher in Tausend},ylabel style={font=\small},tick label style={font=\small},nodes near coords,every node near coord/.append style={font=\scriptsize,/pgf/number format/.cd,use comma,1000 sep={}},ymajorgrids,grid style={mbgitter}]\addplot[fill=mbkasten,draw=black] coordinates {(2007,600.0) (2008,420.0) (2009,320.0) (2010,400.0) (2011,260.0) (2012,260.0) (2013,275.0) (2014,300.0) (2015,250.0)};\end{axis}\end{tikzpicture}\par
\fusshilfe{33: 25 \%, Tipp: prozentuale Zunahme}
\end{pfaufgabe}}
\pfpaar{\pfkopfzeile{34.}{eigene Aufgabe}
Ein Fahrrad kostet $600\,€$. Bei Angebot A gibt es $10\,\%$ Rabatt. Auf den neuen Preis gibt es noch einmal $10\,\%$. Lisa hat $490\,€$ gespart. Reicht ihr Geld für Angebot A?
\fusshilfe{34: Ja}}{}
\pfstufe[saussagenprfen]{Aussagen prüfen}
\kommtdran{in 4 von 13 Jahren (2017, 2019, 2023, 2025) \textperiodcentered{} zuletzt 2025 \textperiodcentered{} 2--3 BE}
\begin{pfaufgabe}{35.}{3 BE \textperiodcentered{} P10 2025 · 4b}
Eine Rampe für Rollstühle darf höchstens 6 \% Steigung haben, damit man sie ohne fremde Hilfe hinauffahren kann. Vervollständigen Sie den Satz: „Bei 6 \% Steigung steigt man auf \_\_\_ Metern waagerechter Strecke um \_\_\_ Meter an.“ Am Hauseingang soll eine Rampe eine 16 cm hohe Stufe überwinden; waagerecht ist sie 170 cm lang. Frau Yücel meint: „Unsere Rampe wird zu steil, ihre Steigung liegt über 6 \%.“ Überprüfen Sie, ob Frau Yücel recht hat.
\par\smallskip\noindent \begin{tikzpicture}[line width=0.6pt]\draw (0,0) -- (6,0) -- (6,1.4) -- cycle;\draw (5.7,0) -- (5.7,0.3) -- (6,0.3);\node[below,font=\small] at (3,0) {170 cm};\node[right,font=\small] at (6,0.7) {16 cm};\node[above left,font=\small] at (3,0.7) {$x$};\node[font=\small] at (1.1,0.12) {$\alpha$};\node[font=\scriptsize,mbgrau,anchor=west] at (0,-0.75) {nicht maßstabsgerecht};\end{tikzpicture}\par
\rechenplatz{3}
\fusshilfe{35: 100 m und 6 m; Ja, Steigung 16 : 170 $\approx$ 9,4 \% > 6 \%, Tipp: Steigung in Prozent}
\end{pfaufgabe}
\pfweiterekopf
\pfpaar{\pfkopfzeile{36.}{eigene Aufgabe nach P10 2020}
$5\,\%$ der Pakete kommen zu spät. Kreuze die Aussage an, die dazu passt.\\ \kreuz{$5$ von $10$ Paketen kommen zu spät.}\\ \kreuz{Jedes 5. Paket kommt zu spät.}\\ \kreuz{$5$ von $100$ Paketen kommen zu spät.}
\fusshilfe{36: $5$ von $100$ Paketen kommen zu spät}}{\pfkopfzeile{37.}{eigene Aufgabe nach P10 2026}
Ein Liter Milch kostete $1{,}15\,€$, jetzt kostet er $1{,}29\,€$. Um wie viel Prozent ist der Preis gestiegen? Runde auf eine Stelle nach dem Komma.
\fusshilfe{37: $\approx 12{,}2\,\%$}}
{\small\begin{pfaufgabe}{38.}{2 BE \textperiodcentered{} P10 2023 · 6b}
Im Jahr 2020 gab es auf der Erde rund 7,8 Milliarden Menschen. Die Tabelle teilt sie nach dem Alter auf: Kinder (0 bis 14 Jahre) 1,95 Milliarden, Jugendliche (15 bis 24 Jahre) 1,25 Milliarden, mittleres Alter (25 bis 64 Jahre) 3,9 Milliarden, höheres Alter (ab 65 Jahre) 0,7 Milliarden. Eine der beiden folgenden Aussagen widerspricht der Tabelle: Kreuzen Sie die falsche Aussage an. Formulieren Sie eine berichtigte Aussage, die zur Tabelle passt.
\pfkreuzzeile{„2020 war jeder zweite Mensch auf der Erde im mittleren Alter.“}\pfkreuzzeile{„2020 war ein Drittel aller Menschen im Kindesalter.“}
\fusshilfe{38: Aussage 2 falsch; Ein Viertel der Weltbevölkerung waren 2020 Kinder (1,95 : 7,8 = 0,25), Tipp: Anteil}
\end{pfaufgabe}}
{\small\begin{pfaufgabe}{39.}{2 BE \textperiodcentered{} P10 2019 · 5b}
Die Tabelle gibt an, welcher Anteil der Mädchen und der Jungen ein Medium in der Freizeit nutzt (Mädchen / Jungen): Smartphone 98 \% / 95 \%, Fernseher 81 \% / 75 \%, Konsolenspiele 14 \% / 72 \%, Bücher 46 \% / 30 \%. Zu der Tabelle gibt es zwei Aussagen: Kreuzen Sie die falsche Aussage an. Berichtigen Sie diese Aussage.
\pfkreuzzeile{„Drei Viertel der Mädchen sehen in ihrer Freizeit fern.“}\pfkreuzzeile{„Bei Jungen sind Konsolenspiele deutlich beliebter als bei Mädchen.“}
\fusshilfe{39: Aussage 1 falsch; etwa vier Fünftel (81 \%) der Mädchen nutzen einen Fernseher, Tipp: Anteil}
\end{pfaufgabe}}
\pfpaar{\pfkopfzeile{40.}{2 BE \textperiodcentered{} P10 2017 · 2b}
Von allen Kindern, die 2013 in Deutschland zur Welt kamen, waren 49,4 \% das erste Kind ihrer Mutter, 34,4 \% das zweite, 11,2 \% das dritte und 5,0 \% das vierte oder ein späteres. Eine Forscherin folgert daraus: „Über die Hälfte dieser Kinder hat mindestens ein Geschwisterkind.“ Geben Sie eine Begründung für diese Folgerung an.
\fusshilfe{40: 34,4 \% + 11,2 \% + 5,0 \% = 50,6 \% > 50 \%, Tipp: Anteile addieren}}{}
\pfstufe[sprozentauseinerberec]{Prozent aus einer berechneten Fläche}
\kommtdran{in 1 von 13 Jahren (2023) \textperiodcentered{} zuletzt 2023 \textperiodcentered{} 4 BE \textperiodcentered{} \textbf{selten geprüft}}
\begin{pfaufgabe}{41.}{4 BE \textperiodcentered{} P10 2023 · 5c \textperiodcentered{} kennst du aus Flächen: Anteil in Prozent (Verschnitt)}
Für Konservendosen werden kreisrunde Deckel mit dem Radius 4 cm aus Blechstreifen ausgestanzt, die 8 cm breit und 1 m lang sind. Aus jedem Streifen sollen möglichst viele Deckel entstehen. Der Abfall darf höchstens 25 \% eines Streifens ausmachen. Überprüfen Sie rechnerisch, ob diese Vorgabe eingehalten wird.
\rechenplatz{4}
\fusshilfe{41: Ja, 12 Deckel, Abfall $\approx$ 196,8 cm² $\approx$ 24,6 \% $\le$ 25 \%, Tipp: Verschnitt}
\end{pfaufgabe}
\pfweiterekopf
\pfpaar{\pfkopfzeile{42.}{eigene Aufgabe nach P10 2023}
Aus einer quadratischen Korkplatte mit $40$ cm Seitenlänge werden kreisrunde Untersetzer mit dem Radius $5$ cm ausgeschnitten, so viele wie möglich, in Reihen nebeneinander. Berechne, wie viel Prozent der Platte Abfall sind. Runde auf eine Stelle nach dem Komma.
\fusshilfe{42: $\approx 21{,}5\,\%$}}{\pfkopfzeile{43.}{eigene Aufgabe}
Ein Bäcker sticht aus quadratischen Teigplatten je eine möglichst große runde Pizza aus. Er sagt: „Egal wie groß die Platte ist – der Abfall ist immer gut $21\,\%$.“ Überprüfe die Aussage an einer Platte mit $30$ cm Seitenlänge. Begründe, warum die Größe der Platte keine Rolle spielt.
\fusshilfe{43: Ja}}
\pfstufe[szinsenundzinssatz]{Zinsen und Zinssatz}
\kommtdran{in 2 von 13 Jahren (2014, 2015) \textperiodcentered{} zuletzt 2015 \textperiodcentered{} 1 BE \textperiodcentered{} \textbf{selten geprüft}}
\begin{pfaufgabe}{44.}{1 BE \textperiodcentered{} P10 2015 · 1e}
Auf einem Konto liegen 400 €, die mit 2 \% im Jahr verzinst werden. Geben Sie an, wie viel Euro Zinsen nach einem Jahr dazukommen.
\rechenplatz{2}
\fusshilfe{44: 8 €, Tipp: Kapital}
\end{pfaufgabe}
\pfweiterekopf
\pfpaar{\pfkopfzeile{45.}{eigene Aufgabe}
900\,€ Kapital bringen in einem Jahr 27\,€ Zinsen. Wie hoch ist der Zinssatz?
\par\noindent\hfill \leerfeld[\%]
\fusshilfe{45: 3\,\%}}{\pfkopfzeile{46.}{1 BE \textperiodcentered{} P10 2014 · 1e}
Auf einem Sparkonto liegen 10 000 €. Nach einem Jahr werden 230 € Zinsen gutgeschrieben. Bestimmen Sie den Zinssatz.
\fusshilfe{46: 2,3 \%, Tipp: Kapital}}
{\small\begin{pfaufgabe}{47.}{1 BE \textperiodcentered{} P10 2014 · 3a}
Lena zahlt seit ihrem 11. Geburtstag jedes Jahr an ihrem Geburtstag 200,00 € auf ein Sparbuch ein. Am 15. Geburtstag hält sie in einer Tabelle fest, wie sich ihr Guthaben entwickelt hat. Das Guthaben wird mit 2 \% pro Jahr verzinst. Weisen Sie dies anhand der Werte für das erste Jahr nach.
\par\smallskip\noindent \sachtabelle{lrrrrr}{Geburtstag & 11. & 12. & 13. & 14. & 15.}{Zinsen & – & 4,00 € & 8,08 € & 12,24 € & \leerzelle\\ Einzahlung & 200,00 € & 200,00 € & 200,00 € & 200,00 € & 200,00 €\\ Guthaben & 200,00 € & 404,00 € & 612,08 € & \leerzelle & 1040,81 €}\par
\fusshilfe{47: 200 € · 0,02 = 4,00 €, Tipp: Kapital}
\end{pfaufgabe}}
\pfpaar{\pfkopfzeile{48.}{eigene Aufgabe nach P10 2015}
Ein Sparbetrag von 2\,600\,€ wird ein Jahr lang mit 2\,\% pro Jahr verzinst. Wie viel Euro Zinsen gibt es nach einem Jahr?
\par\noindent\hfill \leerfeld[€]
\fusshilfe{48: 52\,€}}{}
\pfstufe[szinseszinsundguthabe]{Zinseszins und Guthabentabelle}
\kommtdran{in 1 von 13 Jahren (2014) \textperiodcentered{} zuletzt 2014 \textperiodcentered{} 2 BE \textperiodcentered{} \textbf{selten geprüft}}
\begin{pfaufgabe}{49.}{2 BE \textperiodcentered{} P10 2014 · 3b}
Lena zahlt seit ihrem 11. Geburtstag jedes Jahr an ihrem Geburtstag 200,00 € auf ein Sparbuch ein. Am 15. Geburtstag hält sie in einer Tabelle fest, wie sich ihr Guthaben entwickelt hat. Vervollständigen Sie die fehlenden Einträge der Tabelle.
\par\smallskip\noindent \sachtabelle{lrrrrr}{Geburtstag & 11. & 12. & 13. & 14. & 15.}{Zinsen & – & 4,00 € & 8,08 € & 12,24 € & \leerzelle\\ Einzahlung & 200,00 € & 200,00 € & 200,00 € & 200,00 € & 200,00 €\\ Guthaben & 200,00 € & 404,00 € & 612,08 € & \leerzelle & 1040,81 €}\par
\rechenplatz{2}
\fusshilfe{49: 824,32 €; 16,49 €, Tipp: Zinsen}
\end{pfaufgabe}
\pfweiterekopf
{\small\begin{pfaufgabe}{50.}{eigene Aufgabe nach P10 2014}
Für den Führerschein zahlt Lukas jedes Jahr 350\,€ auf ein Konto ein. Der Zinssatz ist 3\,\%. Vervollständige die Tabelle. Kontrolliere mit dem Guthaben in Jahr 4.
\par\smallskip\noindent \sachtabelle{lrrr}{Jahr & Zinsen & Einzahlung & Guthaben}{1 & – & 350,00\,€ & 350,00\,€\\ 2 & 10,50\,€ & 350,00\,€ & 710,50\,€\\ 3 & 21,32\,€ & 350,00\,€ & \leerzelle\\ 4 & \leerzelle & 350,00\,€ & 1\,464,27\,€}\par
\par\noindent\hfill Guthaben \leerfeld[€], Zinsen \leerfeld[€]\par
\fusshilfe{50: Guthaben = 1\,081,82\,€, Zinsen = 32,45\,€}
\end{pfaufgabe}}
\pfpaar{\pfkopfzeile{51.}{2 BE \textperiodcentered{} P10 2014 · 3c}
Bei einer anderen Bank legt Lena 1 000 € für 5 Jahre an; der Zinssatz liegt bei 3 \% pro Jahr. Berechnen Sie das Guthaben auf diesem Konto nach den 5 Jahren.
\fusshilfe{51: $\approx$ 1159,27 €, Tipp: Wachstumsfaktor}}{\pfkopfzeile{52.}{eigene Aufgabe}
Lara bekommt zur Jugendweihe 2\,000\,€ und legt sie zu 3\,\% mit Zinseszins an. Der Führerschein kostet in 4 Jahren etwa 2\,250\,€. Reicht das Geld für den Führerschein?
\par\noindent\hfill Antwort: \leerfeld
\fusshilfe{52: Ja, knapp: nach 4 Jahren $\approx$ 2\,251,02\,€}}
\label{pruefstein}
\pfpruefstein{P10 2015 · 2}
\begin{pfaufgabe}{53.}{P10 2015 · 2}
Für den Berliner Fernsehturm gelten diese Eintrittspreise: Erwachsene zahlen 12,00 €, Kinder von 3 bis 16 Jahren 7,50 €. In den Wintermonaten gibt es auf den Eintritt 20 \% Rabatt (Rabatt heißt Preisnachlass).
\end{pfaufgabe}
\begin{pfaufgabe}{a)}{3 BE}
Im Winter kommen Vater und Mutter mit ihrem 13-jährigen Sohn und dem Großvater zum Fernsehturm und nutzen den Winterrabatt. Entscheiden Sie für jeden Term, ob er den Rabatt richtig berechnet. Kreuzen Sie richtig oder falsch an:
\par\smallskip\noindent \begingroup\renewcommand{\arraystretch}{1.5}\begin{tabular}{|l|c|c|}\hline Term & richtig & falsch\\ \hline $(2\cdot 12\,€ + 2\cdot 7{,}50\,€)\cdot\tfrac{20}{100}$ & $\square$ & $\square$ \\ $\tfrac{1}{5}\cdot(36\,€ + 7{,}50\,€)$ & $\square$ & $\square$ \\ $\tfrac{20\cdot(12\,€ + 12\,€ + 12\,€ + 7{,}50\,€)}{100}$ & $\square$ & $\square$ \\ \hline\end{tabular}\endgroup\par
\rechenplatz{3}
\end{pfaufgabe}
\begin{pfaufgabe}{b)}{2 BE}
Berechnen Sie den Betrag in Euro, den diese Familie insgesamt für den Eintritt bezahlt.
\rechenplatz{3}
\end{pfaufgabe}
\end{document}